\chapter{Architecture et Base de données}

\section{Introduction}

Ce chapitre détaille l'architecture technique de la plateforme, le choix des technologies, l'organisation modulaire du code et le dictionnaire complet des modèles Django.

\section{Architecture globale}

\subsection{Vue d'ensemble}

L'application repose sur une architecture moderne Full-Stack Client-Serveur à trois niveaux :

\begin{itemize}
\item Frontend : React 18 compilé via Vite, avec Tailwind CSS pour le rendu responsive. L'architecture s'appuie sur des composants réutilisables, React Router v6 et un state management optimisé via Context API.
\item Backend : Django 4.2 LTS et Django REST Framework (DRF), exposant une API REST sécurisée par JWT (SimpleJWT). Le backend est organisé en 5 applications Django indépendantes.
\item Base de données : PostgreSQL 15 géré par l'ORM Django, assurant la persistance des candidatures, configurations et historiques.
\item Temps réel : Django Channels 4.0 avec Redis comme channel layer pour les WebSockets (notifications en temps réel avec barre de progression).
\end{itemize}

\subsection{Architecture REST stateless}

Le frontend communique avec le backend exclusivement via des endpoints REST sécurisés par JWT. Les réponses sont au format JSON. Le pattern est strictement stateless côté serveur, à l'exception des connexions WebSocket persistantes pour les notifications.

\section{Stack technologique}

\begin{table}[htbp]
\centering
\begin{tabular}{|l|l|l|}
\hline
Technologie & Version & Rôle \\
\hline
Django & 4.2 LTS & Framework backend principal \\
\hline
Django REST Framework & 3.17 & Construction de l'API REST \\
\hline
Django Channels & 4.0 & WebSockets temps réel \\
\hline
SimpleJWT & 5.5 & Authentification par tokens JWT \\
\hline
React & 18 & Framework frontend SPA \\
\hline
Tailwind CSS & 3.x & Framework CSS utilitaire \\
\hline
PostgreSQL & 15 & Base de données relationnelle \\
\hline
Redis & 7.x & Channel Layer WebSocket \\
\hline
pdfplumber + pytesseract & — & Extraction OCR \\
\hline
fuzzywuzzy & 0.18 & Correspondance floue des matières \\
\hline
openpyxl & — & Import/Export Excel \\
\hline
\end{tabular}
\caption{Tableau 7}
\label{tab:table_7}
\end{table}

Table 5:  Stack technologique du projet

\section{Organisation modulaire du backend}

Le backend est structuré en 5 applications Django indépendantes, chacune gérant un domaine métier spécifique :

\begin{table}[htbp]
\centering
\begin{tabular}{|l|l|l|}
\hline
Application & Responsabilité & Fichiers clés \\
\hline
users & Authentification, comptes, profils candidats & models.py, serializers.py, views.py \\
\hline
candidatures & Dossiers, documents, notes, services métier & models.py, views.py, services/ \\
\hline
administration & Filières, diplômes, épreuves, historique & models.py, views.py \\
\hline
scoring & Règles de rejet, configuration scoring & models.py, views.py \\
\hline
notifications & Notifications email et WebSocket & models.py, consumers.py, views.py \\
\hline
\end{tabular}
\caption{Tableau 8}
\label{tab:table_8}
\end{table}

Table 6: Applications Django du backend

Le répertoire candidatures/services/ contient la logique métier critique :

\begin{table}[htbp]
\centering
\begin{tabular}{|l|l|l|}
\hline
Service & Fichier & Rôle \\
\hline
Extraction OCR & extraction.py & Extraction des informations depuis les documents scannés \\
\hline
Fuzzy Matching & fuzzy\_matching.py & Correspondance floue des noms de matières \\
\hline
Règles de rejet & regles.py & Application des 7 types de règles de rejet automatique \\
\hline
Scoring & scoring.py & Calcul du score pondéré par filière \\
\hline
Vérification & verification\_document.py & Vérification d'authenticité des documents \\
\hline
Invitation PDF & invitation\_pdf.py & Génération des convocations au format PDF \\
\hline
\end{tabular}
\caption{Tableau 9}
\label{tab:table_9}
\end{table}

Table 7:  Services métier du module candidatures

\section{Dictionnaire des modèles Django}

Le modèle relationnel est structuré autour de 12 entités principales réparties dans les applications Django. Nous détaillons ci-dessous chaque modèle.

\subsection{Application Users}

Le modèle User étend AbstractBaseUser de Django pour personnaliser l'authentification. L'identification se fait par email et CIN (Carte d'Identité Nationale). Trois rôles sont définis : CANDIDAT, RESPONSABLE et ADMIN, implémentant un contrôle d'accès RBAC.

\begin{table}[htbp]
\centering
\begin{tabular}{|l|l|l|}
\hline
Attribut & Type & Contrainte \\
\hline
id & UUID & Clé primaire auto-générée \\
\hline
email & VARCHAR & UNIQUE, NOT NULL \\
\hline
cin & VARCHAR & UNIQUE, NOT NULL \\
\hline
role & ENUM & CANDIDAT / RESPONSABLE / ADMIN \\
\hline
is\_active & BOOLEAN & Défaut : True \\
\hline
password & VARCHAR & Hashé via PBKDF2 \\
\hline
\end{tabular}
\caption{Tableau 10}
\label{tab:table_10}
\end{table}

Le modèle Candidat est relié en OneToOne à User et stocke le profil détaillé :

\begin{table}[htbp]
\centering
\begin{tabular}{|l|l|l|}
\hline
Attribut & Type & Contrainte \\
\hline
id & UUID & Clé primaire \\
\hline
user\_id & FK → User & OneToOne, CASCADE \\
\hline
nom & VARCHAR & NOT NULL \\
\hline
prenom & VARCHAR & NOT NULL \\
\hline
telephone & VARCHAR & Optionnel \\
\hline
date\_naissance & DATE & NOT NULL \\
\hline
\end{tabular}
\caption{Tableau 11}
\label{tab:table_11}
\end{table}

\subsection{Application Candidatures}

Le modèle Dossier est le modèle central de l'application. Il lie un Candidat à une Filière et contient le statut, les scores calculés et les motifs de rejet éventuels.

\begin{table}[htbp]
\centering
\begin{tabular}{|l|l|l|}
\hline
Attribut & Type & Contrainte \\
\hline
id & UUID & Clé primaire \\
\hline
candidat\_id & FK → Candidat & CASCADE \\
\hline
filiere\_id & FK → Filiere & CASCADE \\
\hline
statut & ENUM & 11 valeurs possibles \\
\hline
score & DECIMAL & Score dossier calculé \\
\hline
score\_final & DECIMAL & 40\% dossier + 60\% épreuve \\
\hline
rang\_final & INT & Classement par filière \\
\hline
mention & ENUM & Auto-calculé \\
\hline
moyenne\_generale & DECIMAL & Moyenne des semestres \\
\hline
\end{tabular}
\caption{Tableau 12}
\label{tab:table_12}
\end{table}

Le modèle Document gère les pièces justificatives uploadées par le candidat :

\begin{table}[htbp]
\centering
\begin{tabular}{|l|l|l|}
\hline
Attribut & Type & Contrainte \\
\hline
id & UUID & Clé primaire \\
\hline
dossier\_id & FK → Dossier & CASCADE \\
\hline
type\_doc & ENUM & DIPLOME / RELEVE / CIN / PHOTO \\
\hline
fichier & FILE & Stocké dans media/dossiers/{uuid}/ \\
\hline
qualite\_ocr & DECIMAL & Score de confiance OCR \\
\hline
\end{tabular}
\caption{Tableau 13}
\label{tab:table_13}
\end{table}

Le modèle NoteSemestre remplace l'ancienne approche par NoteMatiere. Il intègre la moyenne par semestre, la session et la mention calculée dynamiquement :

\begin{table}[htbp]
\centering
\begin{tabular}{|l|l|l|}
\hline
Attribut & Type & Contrainte \\
\hline
id & UUID & Clé primaire \\
\hline
dossier\_id & FK → Dossier & CASCADE \\
\hline
semestre & ENUM & S1 à S6 \\
\hline
moyenne & DECIMAL & Plage stricte [10.00, 20.00] \\
\hline
session & ENUM & NORMALE / RATTRAPAGE \\
\hline
mention & ENUM & Auto-calculé (editable=False) \\
\hline
\end{tabular}
\caption{Tableau 14}
\label{tab:table_14}
\end{table}

\subsection{Application Administration et Scoring}

Le modèle Filiere définit les filières du cycle ingénieur avec leurs paramètres :

\begin{table}[htbp]
\centering
\begin{tabular}{|l|l|l|}
\hline
Attribut & Type & Contrainte \\
\hline
id & UUID & Clé primaire \\
\hline
nom & VARCHAR & NOT NULL \\
\hline
code & VARCHAR & UNIQUE (TDI, IACS, IAA, G2ER) \\
\hline
niveau & ENUM & BAC2 / BAC3 \\
\hline
places\_disponibles & INT & Défaut : 30 \\
\hline
\end{tabular}
\caption{Tableau 15}
\label{tab:table_15}
\end{table}

Les modèles EpreuveEcrite et NoteEcrite gèrent la logistique du concours écrit :

\begin{table}[htbp]
\centering
\begin{tabular}{|l|l|l|}
\hline
Modèle & Attributs clés & Rôle \\
\hline
EpreuveEcrite & nom, date\_epreuve, seuil\_admission, statut & Gestion de l'épreuve par filière \\
\hline
NoteEcrite & note, resultat, rang\_final & Note individuelle et classement \\
\hline
ConfigScoring & matiere, poids, bonus\_mention & Pondération scoring par filière \\
\hline
RegleRejet & type\_regle, parametre, is\_active, message & Règle de rejet configurable \\
\hline
HistoriqueAction & action, commentaire, timestamp & Traçabilité des décisions \\
\hline
\end{tabular}
\caption{Tableau 16}
\label{tab:table_16}
\end{table}

Table 8: Modèles d'administration et scoring

\section{Relations et cardinalités}

\begin{table}[htbp]
\centering
\begin{tabular}{|l|l|l|}
\hline
Relation & Type & Cardinalité \\
\hline
User → Candidat & OneToOne & 1..1 \\
\hline
Candidat → Dossier & OneToMany & 1..* \\
\hline
Dossier → Document & OneToMany & 1..4 \\
\hline
Dossier → NoteSemestre & OneToMany & 1..6 \\
\hline
Filiere → Dossier & OneToMany & 1..* \\
\hline
Filiere → RegleRejet & OneToMany & 1..* \\
\hline
Filiere → EpreuveEcrite & OneToMany & 1..* \\
\hline
EpreuveEcrite → NoteEcrite & OneToMany & 1..* \\
\hline
Dossier → HistoriqueAction & OneToMany & 1..* \\
\hline
User → Notification & OneToMany & 1..* \\
\hline
\end{tabular}
\caption{Tableau 17}
\label{tab:table_17}
\end{table}

Table 9: Relations et cardinalités du modèle de données

\section{Machine d'états du dossier}

Le dossier suit une machine d'états stricte implémentée dans la méthode changer\_statut() du modèle. Chaque transition crée automatiquement un enregistrement HistoriqueAction pour la traçabilité :

\begin{table}[htbp]
\centering
\begin{tabular}{|l|l|}
\hline
État & Transitions possibles \\
\hline
BROUILLON & → EN\_TRAITEMENT (soumission) \\
\hline
EN\_TRAITEMENT & → REJETE\_AUTO | EN\_ATTENTE | INCOMPLET | SUSPECT \\
\hline
INCOMPLET & → EN\_TRAITEMENT (complétion) \\
\hline
SUSPECT & → EN\_ATTENTE | REJETE\_FINAL (décision manuelle) \\
\hline
EN\_ATTENTE & → PRESELECTIONNE | REJETE\_FINAL \\
\hline
PRESELECTIONNE & → ADMIS\_FINAL | RECALE\_FINAL | ABSENT\_ECRIT \\
\hline
\end{tabular}
\caption{Tableau 18}
\label{tab:table_18}
\end{table}

Table 10: Tableau 2.6 : Machine d'états du dossier de candidature

\section{Conclusion}

Ce chapitre a détaillé l'architecture Full-Stack de la plateforme et le modèle de données relationnel complet. L'organisation modulaire en 5 applications Django assure une séparation claire des responsabilités. Le chapitre suivant présentera la logique métier et les algorithmes de scoring et d'OCR.

